\chapter{Relevamiento, analisis y definicion del problema}

\section{Necesidad detectada}

Completar con la necesidad comunicacional que motiva el desarrollo del prototipo.

\section{Requisitos funcionales}

\begin{itemize}
    \item Capturar imagen desde la camara del dispositivo movil.
    \item Detectar manos y landmarks en tiempo real.
    \item Construir una secuencia temporal de caracteristicas.
    \item Clasificar la seña realizada.
    \item Mostrar el resultado en pantalla.
    \item Reproducir el resultado mediante sintesis de voz.
\end{itemize}

\section{Requisitos no funcionales}

\begin{itemize}
    \item Ejecutarse localmente en Android.
    \item Reducir la dependencia de conexion a internet.
    \item Mantener tiempos de respuesta aceptables para uso interactivo.
    \item Preservar la privacidad evitando el envio de imagenes a servidores externos.
\end{itemize}

\section{Restricciones del proyecto}

Completar con las restricciones de alcance, datos disponibles, cantidad de señas reconocidas, capacidad de los dispositivos y limitaciones del modelo.

\section{Criterios de evaluacion}

Completar con las metricas consideradas: tiempo de captura, FPS efectivos, tiempo de inferencia, estabilidad de deteccion y comportamiento en distintos dispositivos.
